%%%%%%%%%%%%%%%%%%%%%%%%%%%%%%%%%%%%%%%%%%%%%%%%%%%%%%%%%%%%%%%%%%%%%%%%%%%%%%%%
%% PHASE 2 — DETAILED RESEARCH EXECUTION TABLES
%% Literature Review & Gap Structuring (Chapter 2)
%%%%%%%%%%%%%%%%%%%%%%%%%%%%%%%%%%%%%%%%%%%%%%%%%%%%%%%%%%%%%%%%%%%%%%%%%%%%%%%%
\normalsize\normalfont\color{black}

%% ── Table 1: Input / Process / Output ──
Table~\ref{tab:phase2_ipo} maps the eight steps of the systematic literature review, from keyword strategy through CMM maturity mapping, specifying inputs, processes, and outputs at each stage.

\needspace{4\baselineskip}
\begingroup\singlespacing\tablefontsize\tablestripes
\begin{xltabular}{\textwidth}{@{} l X X X @{}}
\caption[Phase~2 --- Input, Process, and Output: Literature]{Phase~2 --- Input, Process, and Output: Literature Review and Gap Structuring}
\label{tab:phase2_ipo}\\
\toprule
\thead{Step} & \thead{Input} & \thead{Process} & \thead{Output} \\
\midrule
\endfirsthead
\endhead
\bottomrule
\endlastfoot
Keyword strategy       & Research questions (RQ1--RQ8), objectives (O1--O8)               & Define Boolean search strings, MeSH terms, domain-specific keywords & Keyword matrix (6 search strings) \\
\addlinespace
Database search        & Keyword matrix, database access credentials                       & Query IEEE, PubMed, Scopus, Web of Science, Google Scholar          & Raw result set (2,400+ records) \\
\addlinespace
Paper screening        & Raw result set, inclusion/exclusion criteria                      & PRISMA two-stage screening: title/abstract then full-text review    & Filtered corpus (120+ eligible papers) \\
\addlinespace
Theme classification   & Filtered corpus                                                   & Thematic coding across 6 topics: DL, agentic AI, RAG, BCI, governance, trust & 6 thematic clusters with sub-themes \\
\addlinespace
Theory integration     & Thematic clusters, TAM, TOE, RBV, DSR literature                 & Map themes to theoretical lenses; synthesise multi-theory alignment  & Integrated theoretical framework (TAM-TOE-RBV) \\
\addlinespace
Gap identification     & Theoretical framework, thematic clusters                         & Contrast existing studies against research objectives; identify voids & 12 research gaps (G1--G12) \\
\addlinespace
SWOT / Blue Ocean      & Gap analysis, industry context                                   & Assess framework strengths, weaknesses, opportunities, threats; Blue Ocean canvas & SWOT matrix, Blue Ocean strategy canvas \\
\addlinespace
CMM maturity mapping   & Gap analysis, governance literature                              & Classify organisational AI maturity across 5 CMM levels              & AI Governance Maturity Model (Level~1--5) \\
\end{xltabular}
\endgroup

\vspace{1.5em}

%% ── Table 2: Quality Validation Framework ──
Table~\ref{tab:phase2_quality} documents the six quality dimensions that govern the rigour and reproducibility of the literature review process.

\needspace{3\baselineskip}
\begingroup\singlespacing\tablefontsize\tablestripes
\begin{xltabular}{\textwidth}{@{} l X @{}}
\caption{Phase~2 --- Research Design Quality Framework}
\label{tab:phase2_quality}\\
\toprule
\thead{Dimension} & \thead{Method} \\
\midrule
\endfirsthead
\endhead
\bottomrule
\endlastfoot
Quality check        & Ensure systematic, reproducible literature review with transparent search protocol \\
\addlinespace
Justification        & PRISMA-compliant methodology ensures rigour, traceability, and bias minimisation \\
\addlinespace
Framework reference  & PRISMA 2020 guidelines, Kitchenham systematic review protocol, Webster \& Watson concept-centric matrix \\
\addlinespace
Techniques           & Boolean keyword search, forward/backward citation tracing, snowball sampling, thematic coding \\
\addlinespace
Evaluation           & Inter-rater reliability on screening decisions, citation network density, theme saturation check \\
\addlinespace
Benchmarking         & 80--120 papers reviewed, $\geq$3 theoretical lenses, $\geq$4 constructs identified, $\geq$3 research themes \\
\end{xltabular}
\endgroup

\vspace{1.5em}

%% ── Table 3: Academic Readiness Evaluation ──
Table~\ref{tab:phase2_readiness} compares PhD, DBA, and professor-level readiness expectations across six criteria for the literature review and gap structuring phase.

\needspace{4\baselineskip}
\begingroup\singlespacing\tablefontsize\tablestripes
\begin{xltabular}{\textwidth}{@{} l X X X @{}}
\caption{Phase~2 --- Professor, PhD, and DBA Readiness Evaluation}
\label{tab:phase2_readiness}\\
\toprule
\thead{Criteria} & \thead{PhD Readiness} & \thead{DBA Readiness} & \thead{Professor Expectation} \\
\midrule
\endfirsthead
\endhead
\bottomrule
\endlastfoot
Literature coverage    & 120+ papers screened from 2,400 records via PRISMA; 6 Boolean search strings across IEEE, PubMed, Scopus, Web of Science & 70+ applied sources blending WHO/FDA/EU AI Act reports with practitioner studies on clinical AI adoption barriers & Examiner checks PRISMA flow (2,400 $\rightarrow$ 340 $\rightarrow$ 120+); citation recency $\geq$60\% within 5 years confirmed \\
\addlinespace
Theory integration     & TAM-TOE-RBV integrated model with DSR validation lens; each construct (IV, M1, M2, DV, moderator) traced to a theory & TAM explains adoption intent, TOE contextualises organisational readiness, RBV justifies 525-module taxonomy as VRIN resource & Examiner expects explicit mapping: RQ1 $\rightarrow$ TAM, RQ2 $\rightarrow$ TOE, RQ3 $\rightarrow$ RBV, RQ4--6 $\rightarrow$ DSR artefact \\
\addlinespace
Gap articulation       & 12 gaps (G1--G12): fragmented pipelines (G1), no unified governance (G5), limited explainability (G8), trust deficit (G12); H0 null framing & Gaps translate to business need: \$4.5T health burden unaddressed by siloed AI tools; ROI 135--2{,}717\% justifies investment & Each gap traceable to $\geq$3 sources; G1--G12 collectively motivate all 5 hypotheses and 8 research questions \\
\addlinespace
Systematic method      & PRISMA 2020 full protocol: 6 search strings, 5 databases, two-stage screening (title/abstract then full-text), flow diagram with counts & Reproducible strategy with documented inclusion/exclusion criteria; forward/backward citation tracing (2 rounds, saturation check) & Transparent audit trail: keyword matrix published; screening decisions logged; inter-rater reliability on screening decisions \\
\addlinespace
Critical analysis      & Contradictions explored: Heinsfeld (2018) 70\% fMRI-only vs.\ this study's multi-modal 97.67\% ASD; depression heterogeneity acknowledged & Practitioner limitations identified: single-domain tools lack governance; no prior agentic mediation measurement in biomedical AI & Independent critical thinking demonstrated: SWOT analysis, Blue Ocean canvas, CMM maturity model, null hypothesis (H0) framing \\
\addlinespace
Thematic organisation  & 6 themes: (1) DL classification, (2) agentic automation, (3) RAG explainability, (4) BCI integration, (5) AI governance, (6) clinical trust & Themes map directly to RGAIG 5-layer architecture and 525-module taxonomy tiers (DL, GenAI, governance) & Webster \& Watson concept-centric matrix organises themes by construct rather than chronology; supports RQ1--RQ8 alignment \\
\end{xltabular}
\endgroup

\vspace{1.5em}

%% ── Table 4: Research Methodology Consideration ──
Table~\ref{tab:phase2_methodology} summarises the three methodological approaches evaluated for the literature review, identifying the adopted mixed-methods strategy.

\needspace{3\baselineskip}
\begingroup\singlespacing\tablefontsize\tablestripes
\begin{xltabular}{\textwidth}{@{} l l X @{}}
\caption{Phase~2 --- Research Methodology Type Consideration}
\label{tab:phase2_methodology}\\
\toprule
\thead{Method} & \thead{This Study} & \thead{Rationale} \\
\midrule
\endfirsthead
\endhead
\bottomrule
\endlastfoot
Quantitative       & Primary    & Citation frequency analysis, bibliometric mapping, publication trend statistics \\
\addlinespace
Qualitative        & Secondary  & Thematic coding of literature themes, narrative synthesis of theoretical lenses \\
\addlinespace
Mixed methods      & Adopted    & Quantitative bibliometric data triangulated with qualitative thematic synthesis \\
\end{xltabular}
\endgroup

\vspace{1.5em}

%% ── Table 5: Variable Taxonomy (Phase 2 Discovery) ──
Table~\ref{tab:phase2_variables} enumerates the six research variables discovered during the literature synthesis, classifying each by type and theoretical provenance.

\needspace{4\baselineskip}
\begingroup\singlespacing\tablefontsize\tablestripes
\begin{xltabular}{\textwidth}{@{} l l X @{}}
\caption{Phase~2 --- Research Variable Taxonomy Discovered from Literature}
\label{tab:phase2_variables}\\
\toprule
\thead{Variable} & \thead{Type} & \thead{Definition / Literature Source} \\
\midrule
\endfirsthead
\endhead
\bottomrule
\endlastfoot
AI governance framework         & Independent variable  & Organisational policies, standards, and controls governing AI deployment \citep{jobin2019ethics} \\
\addlinespace
AI trust                        & Mediator              & Stakeholder confidence in AI system reliability, fairness, and transparency \\
\addlinespace
AI adoption                     & Dependent variable    & Organisational uptake and sustained use of AI-driven diagnostic systems \\
\addlinespace
Digital culture                 & Moderator             & Organisational readiness, data literacy, and innovation orientation \\
\addlinespace
Regulatory compliance           & Control variable      & External regulatory requirements (FDA, EU AI Act, HIPAA) held constant \\
\addlinespace
Dataset characteristics         & Control variable      & Domain type, sample size, signal modality (EEG, MRI, clinical) \\
\end{xltabular}
\endgroup

\vspace{1.5em}

%% ── Table 6: Systematic Literature Search Strategy ──
Table~\ref{tab:phase2_search} reports the six Boolean search strings, target databases, initial result counts, and post-screening yields that produced the filtered corpus.

\needspace{4\baselineskip}
\begingroup\singlespacing\tablefontsize\tablestripes
\begin{xltabular}{\textwidth}{@{} X l c c @{}}
\caption{Phase~2 --- Systematic Literature Search Strategy}
\label{tab:phase2_search}\\
\toprule
\thead{Search String} & \thead{Databases} & \thead{Results} & \thead{After Screening} \\
\midrule
\endfirsthead
\endhead
\bottomrule
\endlastfoot
``AI governance'' AND ``healthcare'' AND ``diagnosis''                              & Scopus, WoS    & 342 & 48 \\
\addlinespace
``deep learning'' AND ``biomedical'' AND ``ASD-focused''                        & IEEE, PubMed   & 518 & 62 \\
\addlinespace
``responsible AI'' AND ``clinical'' AND ``trust''                                   & Scopus, ACM    & 287 & 41 \\
\addlinespace
``agentic AI'' OR ``AI pipeline'' AND ``automation''                                & IEEE, arXiv    & 196 & 28 \\
\addlinespace
``RAG'' OR ``retrieval augmented'' AND ``explainability''                           & ACM, Scopus    & 214 & 35 \\
\addlinespace
``deep learning'' AND ``EEG'' AND (``autism'' OR ``Parkinson'' OR ``epilepsy'')     & PubMed, IEEE   & 843 & 67 \\
\midrule
\textbf{Total}                                                                      &                & \textbf{2,400} & \textbf{281} \\
\end{xltabular}
\par\vspace{2pt}
\noindent\begin{minipage}{\textwidth}\textit{Note: Search conducted 2023--2025. Duplicates removed across databases. 2,400 initial results $\rightarrow$ 340 after title/abstract screening $\rightarrow$ 120+ after full-text review. PRISMA flow diagram in Chapter~2, Section~2.3.}\end{minipage}
\endgroup

\vspace{1.5em}

%% ── Table 7: Academic Database Selection ──
Table~\ref{tab:phase2_databases} catalogues the six academic databases selected for the systematic review, together with their coverage scope and relevance to this study.

\needspace{4\baselineskip}
\begingroup\singlespacing\tablefontsize\tablestripes
\begin{xltabular}{\textwidth}{@{} l X X @{}}
\caption{Phase~2 --- Academic Database Selection and Justification}
\label{tab:phase2_databases}\\
\toprule
\thead{Database} & \thead{Coverage} & \thead{Relevance to Study} \\
\midrule
\endfirsthead
\endhead
\bottomrule
\endlastfoot
Scopus                  & 27,000+ journals; broad multidisciplinary coverage             & AI governance, information systems, healthcare systems literature \\
\addlinespace
Web of Science          & 21,000+ journals; high-impact only                             & TAM, TOE, RBV theoretical foundations; citation tracking \\
\addlinespace
IEEE Xplore             & 5M+ technical documents                                        & DL architectures, EEG signal processing, BCI systems \\
\addlinespace
PubMed                  & 36M+ biomedical citations                                      & Clinical diagnostics, disease-specific studies, FDA/SaMD \\
\addlinespace
ACM Digital Library     & Computing and information systems research                     & RAG, agentic AI, GenAI quality assurance \\
\addlinespace
Google Scholar          & Broadest coverage; grey literature                             & Working papers, pre-prints, industry reports (WHO, EU AI Act) \\
\end{xltabular}
\endgroup

\vspace{1.5em}

%% ── Table 8: Literature Thematic Clustering ──
Table~\ref{tab:phase2_themes} maps the six thematic clusters identified in the literature to their key sources, principal findings, and the specific research gaps each cluster reveals.

\needspace{4\baselineskip}
\begingroup\singlespacing\tablefontsize\tablestripes
\begin{xltabular}{\textwidth}{@{} p{0.15\textwidth} X X X @{}}
\caption{Phase~2 --- Thematic Clustering of Literature}
\label{tab:phase2_themes}\\
\toprule
\thead{Theme} & \thead{Key Sources} & \thead{Key Findings} & \thead{Gap Identified} \\
\midrule
\endfirsthead
\endhead
\bottomrule
\endlastfoot
DL diagnostics          & Rajpurkar (2017), Heinsfeld (2018), Acharya (2018)    & Domain-specific models achieve 85--98\% accuracy; no ASD-focused framework              & G1: No unified ASD-focused architecture \\
\addlinespace
Agentic AI automation      & Lewis (2023), Park (2023), \citet{wang2024agentai}                & Pipeline orchestration improves throughput; limited healthcare validation                 & G5: No mediation evidence in biomedical AI \\
\addlinespace
RAG explainability         & Lewis et al.\ (2020), Gao (2024)                     & RAG improves LLM factuality; no domain-specific clinical RAG evaluation                  & G6: No clinical trust validation for RAG \\
\addlinespace
AI governance frameworks   & Jobin (2019), Floridi (2019), EU AI Act (2024)        & 84 governance frameworks mapped; none integrate DL + GenAI + governance                  & G7--G12: No 525-module unified taxonomy \\
\addlinespace
Clinical trust \& adoption & \citet{davis1989tam}, \citet{venkatesh2003utaut}, Cai (2019)            & TAM explains adoption; trust is mediating factor; limited biomedical AI evidence          & G2: No TAM-TOE-RBV integration for AI diagnostics \\
\addlinespace
BCI \& wearable sensing    & Mullen (2015), Craik (2019), Ahn (2022)               & EEG-based classification feasible; real-time BCI underexplored                           & G4: No wearable BCI integration with governance \\
\end{xltabular}
\endgroup

\vspace{1.5em}

%% ── Table 9: Evaluation and Benchmarking ──
Table~\ref{tab:phase2_benchmark} enumerates the eight benchmarking standards against which the literature review quality is assessed.

\needspace{4\baselineskip}
\begingroup\singlespacing\tablefontsize\tablestripes
\begin{xltabular}{\textwidth}{@{} l X @{}}
\caption{Phase~2 --- Evaluation Metrics and Benchmarking}
\label{tab:phase2_benchmark}\\
\toprule
\thead{Metric / Benchmark} & \thead{Standard} \\
\midrule
\endfirsthead
\endhead
\bottomrule
\endlastfoot
Literature volume             & 80--120 peer-reviewed papers screened through PRISMA protocol \\
\addlinespace
Theoretical lenses            & 3--5 complementary theories (TAM, TOE, RBV, DSR, Diffusion of Innovation) \\
\addlinespace
Construct identification      & 4--6 constructs with clear operationalisation pathway \\
\addlinespace
Thematic coverage             & 3--5 major themes with sub-theme decomposition \\
\addlinespace
Gap completeness              & All 12 gaps (G1--G12) traceable to at least 3 supporting sources \\
\addlinespace
Citation recency              & $\geq$60\% of sources published within the last 5 years \\
\addlinespace
PRISMA compliance             & Full PRISMA 2020 flow diagram with identification, screening, eligibility, inclusion counts \\
\addlinespace
Forward/backward tracing      & At least 2 rounds of citation snowballing to ensure saturation \\
\end{xltabular}
\endgroup

\vspace{1.5em}

%% ── Table 10: Research Evidence Chain ──
Table~\ref{tab:phase2_evidence} reports the eight-link evidence chain tracing the literature review from keyword strategy through to the CMM maturity model.

\needspace{4\baselineskip}
\begingroup\singlespacing\tablefontsize\tablestripes
\begin{xltabular}{\textwidth}{@{} l X @{}}
\caption{Phase~2 --- Research Evidence Chain}
\label{tab:phase2_evidence}\\
\toprule
\thead{Chain Link} & \thead{Evidence} \\
\midrule
\endfirsthead
\endhead
\bottomrule
\endlastfoot
Keyword strategy             & 6 Boolean search strings covering DL diagnostics, agentic AI, RAG, BCI, AI governance, clinical trust \\
\addlinespace
Database results             & 2,400+ records retrieved from IEEE, PubMed, Scopus, Web of Science, Google Scholar \\
\addlinespace
PRISMA screening             & Two-stage screening: 2,400 $\rightarrow$ 340 (title/abstract) $\rightarrow$ 120+ (full-text eligibility) \\
\addlinespace
Literature classification    & 6 thematic topics: (1) DL classification, (2) agentic automation, (3) RAG explainability, (4) BCI integration, (5) AI governance, (6) clinical trust \\
\addlinespace
Theoretical framework        & TAM-TOE-RBV integrated model with DSR validation lens \\
\addlinespace
Research gaps                & 12 gaps (G1--G12): fragmented pipelines, no unified governance, limited explainability, trust deficit \\
\addlinespace
SWOT analysis                & Strengths: ASD-focused scope; Weaknesses: data heterogeneity; Opportunities: regulatory alignment; Threats: rapid model obsolescence \\
\addlinespace
CMM maturity model           & 5-level AI governance maturity: Ad Hoc $\rightarrow$ Repeatable $\rightarrow$ Defined $\rightarrow$ Managed $\rightarrow$ Optimising \\
\end{xltabular}
\endgroup

\vspace{1.5em}

%% ── Table 11: Quality Checklist ──
Table~\ref{tab:phase2_checklist} presents the ten-item quality checklist with pass/fail status and supporting evidence for the literature review phase.

\needspace{3\baselineskip}
\begingroup\singlespacing\tablefontsize\tablestripes
\begin{xltabular}{\textwidth}{@{} p{0.33\textwidth} c X @{}}
\caption{Phase~2 --- Quality Checklist}
\label{tab:phase2_checklist}\\
\toprule
\thead{Checklist Item} & \thead{Status} & \thead{Evidence} \\
\midrule
\endfirsthead
\endhead
\bottomrule
\endlastfoot
Is the literature search strategy documented and reproducible?      & $\checkmark$ & Systematic review protocol with 6 Boolean search strings in Section~2.3 \\
\addlinespace
Is the PRISMA flow diagram complete with all stage counts?          & $\checkmark$ & PRISMA 2020 flow diagram: 2,400 $\rightarrow$ 340 $\rightarrow$ 120+ papers (Section~2.4) \\
\addlinespace
Are inclusion/exclusion criteria explicitly stated?                 & $\checkmark$ & Inclusion/exclusion table in Section~2.3 with date range, language, and domain filters \\
\addlinespace
Are at least 80 peer-reviewed sources included?                     & $\checkmark$ & 217 sources in references.bib; 120+ eligible after full-text screening \\
\addlinespace
Are theoretical lenses identified and justified?                    & $\checkmark$ & TAM-TOE-RBV integration justified in Section~2.5; DSR lens in Section~2.6 \\
\addlinespace
Are research gaps traceable to literature evidence?                 & $\checkmark$ & 12 gaps (G1--G12) each mapped to $\geq$3 supporting sources in Section~2.7 \\
\addlinespace
Is forward/backward citation tracing conducted?                     & $\checkmark$ & 2 rounds of snowball sampling documented in Section~2.3 search protocol \\
\addlinespace
Are themes organised using a concept-centric approach?              & $\checkmark$ & Webster \& Watson concept-centric matrix across 6 thematic topics in Section~2.4 \\
\addlinespace
Is SWOT analysis aligned with identified gaps?                      & $\checkmark$ & SWOT matrix in Section~2.8 cross-referenced to G1--G12; Blue Ocean canvas in Section~2.9 \\
\addlinespace
Is the CMM maturity model grounded in governance literature?        & $\checkmark$ & 5-level AI Governance Maturity Model in Section~2.10 citing SEI-CMM and ISO/IEC 38500 \\
\end{xltabular}
\endgroup

\vspace{1.5em}

%% ═══════════════════════════════════════════════════════════════════════════════
%% RGAIG+ THEORETICAL RESEARCH MODEL — Integrated from rgaig_theory_tables.tex
%% Phase 2 (Ch.2): Literature Synthesis + Gap-RQ-Hypothesis Matrix
%% ═══════════════════════════════════════════════════════════════════════════════

%% ── Table 12: Literature Synthesis (RGAIG+ Theory Table 6) ──
Table~\ref{tab:rgaig_lit_synthesis} synthesises the eight literature themes, linking each to its key findings, the identified research gap, and the corresponding RGAIG+ framework design implication.

\needspace{4\baselineskip}
\begingroup\singlespacing\tablefontsize\tablestripes
\begin{xltabular}{\textwidth}{@{} l X X X @{}}
\caption[Phase~2 --- Literature Synthesis: Themes, Gaps, and]{Phase~2 --- Literature Synthesis: Themes, Gaps, and RGAIG+ Framework Implications}
\label{tab:rgaig_lit_synthesis}\\
\toprule
\thead{Theme} & \thead{Key Findings} & \thead{Research Gap} & \thead{RGAIG+ Implication} \\
\midrule
\endfirsthead
\endhead
\bottomrule
\endlastfoot
AI Governance & NIST AI RMF, ISO/IEC 42001, EU AI Act provide domain-general guidance & No consumer wearable-specific governance exists & 9-layer architecture with device-aware governance \\
\addlinespace
Consumer EEG & Consumer devices (4--16 ch) achieve 85--97\% accuracy in controlled settings & Signal quality not integrated with governance & Device Layer (L1) with SNR monitoring \\
\addlinespace
GenAI in Healthcare & LLM-based systems improve diagnostic explanation & Hallucination risk unaddressed for consumer AI & GenAI Engine (L4) with hallucination detection \\
\addlinespace
Trust in AI & Trust in Automation: performance, process, purpose dimensions & Trust undefined for consumer AI diagnostics & Trust mediator with 3-dimensional instrument \\
\addlinespace
Technology Adoption & TAM: perceived usefulness and ease of use drive adoption & TAM not applied to governed AI diagnostics & Extended TAM with governance antecedents \\
\addlinespace
Agentic AI & Multi-agent orchestration improves task coordination & No governance for agentic AI in clinical use & Agentic Orchestration Layer (L5) \\
\addlinespace
Responsible AI & Fairness, accountability, transparency well-established & Implementation gap: principles to controls & Cross-cutting pillar across all 9 layers \\
\addlinespace
Safety/Regulation & IEC 60601, ISO 14971, FDA SaMD provide foundations & Consumer devices outside traditional scope & Safety Module (L8) adapted for consumers \\
\end{xltabular}
\endgroup

\vspace{1.5em}

%% ── Table 13: Gap → RQ → Hypothesis → Framework Component Matrix (RGAIG+ Theory Table 5) ──
Table~\ref{tab:rgaig_gap_rq_hypothesis} maps each identified research gap (G1--G6) to its corresponding research question, structural hypothesis, RGAIG+ framework component, and validation method.

\needspace{4\baselineskip}
\begingroup\singlespacing\tablefontsize\tablestripes
\begin{xltabular}{\textwidth}{@{} c p{0.14\textwidth} X l p{0.14\textwidth} l @{}}
\caption[Phase~2 --- Research Alignment Matrix: Gap → Research]{Phase~2 --- Research Alignment Matrix: Gap → Research Question → Hypothesis → RGAIG+ Component}
\label{tab:rgaig_gap_rq_hypothesis}\\
\toprule
\thead{ID} & \thead{Research Gap} & \thead{Research Question} & \thead{Hyp.} & \thead{RGAIG+ Component} & \thead{Validation} \\
\midrule
\endfirsthead
\endhead
\bottomrule
\endlastfoot
G1 & No unified governance for consumer AI diagnostics & How can governance controls improve trust? & SH1 & Governance (L9) & Expert, SEM \\
\addlinespace
G2 & Lack of XAI standards for consumer GenAI & Role of AI explainability in trust? & SH2 & GenAI Engine (L4) & User study \\
\addlinespace
G3 & Insufficient safety assurance for wearable AI & How does safety affect trust? & SH3 & Safety Module (L8) & Computational \\
\addlinespace
G4 & Device reliability not governed & Does reliability influence trust? & SH4 & Device Layer (L1) & Lab testing \\
\addlinespace
G5 & Opaque monitoring undermines confidence & Monitoring transparency → trust? & SH5 & Monitoring (L8) & Survey \\
\addlinespace
G6 & Trust--adoption pathway undefined & Does trust mediate governance → adoption? & SH6--SH7 & Interface (L6) & SEM mediation \\
\end{xltabular}
\par\smallskip\noindent\begin{minipage}{\dimexpr\textwidth-2\tabcolsep}\textit{Note.} G1--G6 = research gaps from Chapter~2. SH1--SH7 = structural hypotheses (Table~\ref{tab:rgaig_sem_hypotheses}). L1--L9 = RGAIG+ layers.\end{minipage}
\endgroup

\clearpage
